\subsection{Examples of Brauer Groups}

The Brauer group \(\mathrm{Br}(\mathrm{K})\) is reduced to the identity element in the following three cases:

\begin{enumerate}
\def\labelenumi{\alph{enumi})}
\tightlist
\item
  K is separably closed (VIII, p.~253, Corollary 1).
\end{enumerate}

*b) K is a finite field (VIII, p.~357, Corollary 2).

\begin{enumerate}
\def\labelenumi{\alph{enumi})}
\setcounter{enumi}{2}
\tightlist
\item
  K has property \((\mathrm{C}_1)\) (VIII, p.~357, Remark 2).
\end{enumerate}

Suppose that K is a maximal ordered field (VI, §2, No.~5, p.~25). The Brauer group of K is then cyclic of order 2; its elements are the class of K and the class of the quaternion K-algebra of type \((-1,0,-1)\) (III, §2, No.~5, p.~444 and VIII, p.~367, Theorem 1).

Suppose that K is a locally compact, nondiscrete, commutative topological field. If K is not connected, then it is a complete field for a discrete valuation, with finite residue field (Comm. Alg., VI, §9, No.~3, p.~433, Theorem 1), and there exists an isomorphism from Br(K) to Q/Z (VIII, p.~332, Exercise 17). If K is connected, then it is isomorphic to R or C. The Brauer group of the field R is cyclic of order 2. Its nontrivial element is the class of the algebra H of Hamilton quaternions (Gen.~Top., VIII, §1, No.~4, p.~104); the Brauer group of C has order 1.

